% Created 2026-05-24 Sun 23:35
% Intended LaTeX compiler: lualatex
\documentclass[11pt]{article}
\usepackage{amsmath}
\usepackage{fontspec}
\usepackage{graphicx}
\usepackage{longtable}
\usepackage{wrapfig}
\usepackage{rotating}
\usepackage[normalem]{ulem}
\usepackage{capt-of}
\usepackage{hyperref}
\usepackage{minted}
\usepackage{fontspec}
\setmainfont{UT Sans}[
BoldFont={UT Sans Bold},
UprightFont={UT Sans Regular},
FontFace={m}{n}{UT Sans Medium}
]
% Fallback fake italic using a slant
\usepackage{etoolbox}
\newcommand{\fakeitalic}[1]{{\addfontfeatures{FakeSlant=0.18}#1}}
\AtBeginDocument{\renewcommand{\emph}[1]{\fakeitalic{#1}}}
\usepackage{minted}
\usepackage[a4paper,margin=3cm]{geometry}
\usepackage{lastpage}
\usepackage{fancyhdr}
\usepackage{lastpage}
\pagestyle{fancy}                % set fancy style
\fancyhf{}                        % clear all headers and footers
\cfoot{\thepage/\pageref*{LastPage}}  % center footer
\renewcommand{\headrulewidth}{0pt}    % remove header line
\renewcommand{\footrulewidth}{0pt}    % remove footer line
% also override plain page style for first page
\fancypagestyle{plain}{%
\fancyhf{}%
\cfoot{\thepage/\pageref*{LastPage}}%
\renewcommand{\headrulewidth}{0pt}%
\renewcommand{\footrulewidth}{0pt}%
}
\usepackage{url}
\author{BAJCSI Elias-Robert}
\date{}
\title{N17 - teme lab07}
\hypersetup{
 pdfauthor={BAJCSI Elias-Robert},
 pdftitle={N17 - teme lab07},
 pdfkeywords={},
 pdfsubject={},
 pdfcreator={Emacs 30.2 (Org mode 9.7.39)}, 
 pdflang={English}}
\begin{document}

\maketitle
\section{Programare awk}
\label{sec:orgb41e9c2}

Problemele de mai jos se vor rezolva prin script-uri shell ce vor apela utilitarul awk
\subsection{Sa se afiseze ora curenta sub forma}
\label{sec:org5a6c7f4}

ora xx, xx minute, xx secunde
\begin{minted}[breaklines=true,breakanywhere=true,breaksymbol=,fontsize=\small,linenos=false]{sh}
#!/usr/bin/env sh
# Afiseaza ora curenta in formatul cerut.
usage() {
    echo "Usage: ${0}" >&2
    echo "Prints current time as: ora xx, xx minute, xx secunde." >&2
}

[ "$#" -ne 0 ] && { usage; exit 1; }

date '+%H %M %S' | awk '{printf "ora %s, %s minute, %s secunde\n", $1, $2, $3}'
\end{minted}

\begin{center}
\includegraphics[width=.9\linewidth]{./screenshots/awk/1.png}
\end{center}
\subsection{Sa se afiseze procesele ce au consumat mai mult de 10s de unitate centrala.}
\label{sec:orgd3d8b86}
\begin{minted}[breaklines=true,breakanywhere=true,breaksymbol=,fontsize=\small,linenos=false]{sh}
#!/usr/bin/env sh
# Afiseaza procesele cu timp CPU total > 10 secunde.
usage() {
    echo "Usage: ${0}" >&2
    echo "Lists processes with total CPU time greater than 10 seconds." >&2
}

[ "$#" -ne 0 ] && { usage; exit 1; }

ps -e -o pid= -o user= -o comm= -o time= | awk '
function to_seconds(t, n,a,s) {
    n = split(t, a, ":")
    if (n == 2) {
        return a[1] * 60 + a[2]
    }
    if (n == 3) {
        return a[1] * 3600 + a[2] * 60 + a[3]
    }
    return 0
}
{
    sec = to_seconds($NF)
    if (sec > 10) {
        printf "PID=%s USER=%s CMD=%s CPU=%s (%ds)\n", $1, $2, $3, $NF, sec
    }
}'
\end{minted}

\begin{center}
\includegraphics[width=.9\linewidth]{./screenshots/awk/2.png}
\end{center}
\subsection{Sa se listeze cele mai mari 10 fisiere dintr-un director dat ca parametru si din subdirectoarele sale, in ordine descrescatoare a lungimii}
\label{sec:org8aa3839}
\begin{minted}[breaklines=true,breakanywhere=true,breaksymbol=,fontsize=\small,linenos=false]{sh}
#!/usr/bin/env sh
# Listeaza cele mai mari 10 fisiere dintr-un director dat.
usage() {
    echo "Usage: ${0} DIR" >&2
    echo "Lists top 10 largest files recursively in DIR." >&2
}

[ "$#" -ne 1 ] && { usage; exit 1; }

find "$1" -type f -exec ls -ln {} + 2>/dev/null |
awk '{print $5, $9}' |
sort -nr |
awk 'NR <= 10 {printf "%s bytes\t%s\n", $1, $2}'
\end{minted}

\begin{center}
\includegraphics[width=.9\linewidth]{./screenshots/awk/3.png}
\end{center}
\subsection{Se da, in fisierul studenti.lst lista studentilor in formatul:}
\label{sec:org530af15}

Nume Prenume nume-utilizator grupa

(coloanele sunt separate prin blancuri; numele compuse sunt separate prin '-'; exemplu: Popescu Ana-Maria pa12345 221 . Dandu-se o lista de grupe, sa se trimita fiecarui student cate un mesaj (mesajul se da intr-un fisier al carui nume este de asemenea dat ca parametru).
\begin{minted}[breaklines=true,breakanywhere=true,breaksymbol=,fontsize=\small,linenos=false]{sh}
#!/usr/bin/env sh
# Trimite mesaj studentilor din grupele cerute pe baza fisierului studenti.lst.
usage() {
    echo "Usage: ${0} STUDENTI_LST MESAJ_TXT GRUPA1 [GRUPA2 ...]" >&2
    echo "Sends message to users from requested groups." >&2
}

[ "$#" -lt 3 ] && { usage; exit 1; }

students_file=$1
message_file=$2
shift 2

if [ ! -f "$students_file" ] || [ ! -f "$message_file" ]; then
    echo "Eroare: fisier inexistent" >&2
    exit 1
fi

awk -v groups="$*" '
BEGIN {
    n = split(groups, g, " ")
    for (i = 1; i <= n; i++) wanted[g[i]] = 1
}
{
    # Format: Nume Prenume user grupa
    if (NF >= 4 && ($4 in wanted)) {
        print $3
    }
}' "$students_file" |
while IFS= read -r user; do
    if command -v mail >/dev/null 2>&1; then
        mail -s "Mesaj laborator" "$user" < "$message_file"
        echo "Trimis catre: $user"
    else
        echo "mail indisponibil; destinatar: $user"
    fi
done
\end{minted}

\begin{center}
\includegraphics[width=.9\linewidth]{./screenshots/awk/4.png}
\end{center}
\subsection{Pornind de la un director, sa se afiseze toate fisierele, din el si din subdirectoarele sale, care au proprietar diferit de proprietarul directorului in care se gasesc.}
\label{sec:org64d15c8}
\begin{minted}[breaklines=true,breakanywhere=true,breaksymbol=,fontsize=\small,linenos=false]{sh}
#!/usr/bin/env sh
# Afiseaza fisierele al caror proprietar difera de proprietarul directorului parinte.
usage() {
    echo "Usage: ${0} DIR" >&2
    echo "Lists files whose owner differs from parent directory owner." >&2
}

[ "$#" -ne 1 ] && { usage; exit 1; }

find "$1" -type f 2>/dev/null | while IFS= read -r f; do
    d=$(dirname "$f")
    f_owner=$(ls -ldn "$f" | awk '{print $3}')
    d_owner=$(ls -ldn "$d" | awk '{print $3}')
    if [ "$f_owner" != "$d_owner" ]; then
        printf "%s (owner file=%s, owner dir=%s)\n" "$f" "$f_owner" "$d_owner"
    fi
done
\end{minted}


\begin{center}
\includegraphics[width=.9\linewidth]{./screenshots/awk/5.png}
\end{center}
\section{Probleme de rezolvat folosind comanda grep}
\label{sec:orgcc089d4}

\subsection{Sa se afiseze continutul tuturor fisierelor text din directorul dat ca parametru in linia de comanda si din subdirectoarele lui.}
\label{sec:org96fbdec}
\begin{minted}[breaklines=true,breakanywhere=true,breaksymbol=,fontsize=\small,linenos=false]{sh}
#!/usr/bin/env sh
# Afiseaza continutul tuturor fisierelor text din director si subdirectoare.
usage() {
    echo "Usage: ${0} DIR" >&2
    echo "Prints content of text files found recursively in DIR." >&2
}

[ "$#" -ne 1 ] && { usage; exit 1; }

find "$1" -type f 2>/dev/null | while IFS= read -r f; do
    if LC_ALL=C grep -q '[^[:print:][:space:]]' "$f" 2>/dev/null; then
        :
    else
        echo "===== $f ====="
        cat "$f"
    fi
done
\end{minted}


\begin{center}
\includegraphics[width=.9\linewidth]{./screenshots/grep/1.png}
\end{center}
\subsection{Sa se afiseze numele fisierelor din linia de comanda care contin un anumit cuvant dat si el ca parametru in linia de comanda. Se cere sa se afiseze si numarul lor.}
\label{sec:org24c90f5}
\begin{minted}[breaklines=true,breakanywhere=true,breaksymbol=,fontsize=\small,linenos=false]{sh}
#!/usr/bin/env sh
# Afiseaza fisierele din linia de comanda care contin cuvantul dat + numarul lor.
usage() {
    echo "Usage: ${0} WORD FILE1 [FILE2 ...]" >&2
    echo "Prints file names containing WORD and their count." >&2
}

[ "$#" -lt 2 ] && { usage; exit 1; }

word=$1
shift

matches=$(grep -l -- "$word" "$@" 2>/dev/null || true)
[ -n "$matches" ] && printf "%s\n" "$matches"
count=$(printf "%s\n" "$matches" | awk 'NF{c++} END{print c+0}')
echo "Numar fisiere: $count"
\end{minted}

\begin{center}
\includegraphics[width=.9\linewidth]{./screenshots/grep/2.png}
\end{center}
\subsection{Sa se afiseze numele tuturor fisierelor binare din directoarele date ca parametri in linia de comanda si din subdirectoarele acestora.}
\label{sec:org7d98c92}
\begin{minted}[breaklines=true,breakanywhere=true,breaksymbol=,fontsize=\small,linenos=false]{sh}
#!/usr/bin/env sh
# Afiseaza numele fisierelor binare din directoarele date (recursiv).
usage() {
    echo "Usage: ${0} DIR1 [DIR2 ...]" >&2
    echo "Prints binary files found recursively in given directories." >&2
}

[ "$#" -lt 1 ] && { usage; exit 1; }

for d in "$@"; do
    find "$d" -type f 2>/dev/null | while IFS= read -r f; do
        if LC_ALL=C grep -q '[^[:print:][:space:]]' "$f" 2>/dev/null; then
            echo "$f"
        fi
    done
done
\end{minted}

\begin{center}
\includegraphics[width=.9\linewidth]{./screenshots/grep/3.png}
\end{center}
\subsection{Sa se afiseze numele tuturor serverelor la care este conectat fiecare utilizator dat ca parametru in linia de comanda.}
\label{sec:orga20b8c6}
\begin{minted}[breaklines=true,breakanywhere=true,breaksymbol=,fontsize=\small,linenos=false]{sh}
#!/usr/bin/env sh
# Afiseaza serverele (host-uri) la care este conectat fiecare utilizator dat.
usage() {
    echo "Usage: ${0} USER1 [USER2 ...]" >&2
    echo "Prints remote servers each user is connected to." >&2
}

[ "$#" -lt 1 ] && { usage; exit 1; }

for u in "$@"; do
    echo "Utilizator: $u"
    who | awk -v usr="$u" '
        $1 == usr {
            host = $NF
            gsub(/[()]/, "", host)
            if (host != "") seen[host] = 1
        }
        END {
            n = 0
            for (h in seen) {
                print "  " h
                n++
            }
            if (n == 0) print "  (niciun server remote detectat)"
        }'
done
\end{minted}

\begin{center}
\includegraphics[width=.9\linewidth]{./screenshots/grep/4.png}
\end{center}
\subsection{Sa se afiseze numele tuturor utilizatorilor din sistem care fac parte dintr-un anumit grup al carui identificator este dat ca parametru in linia de comanda. Sa se afiseze si numele grupului.}
\label{sec:orge350b13}
\begin{minted}[breaklines=true,breakanywhere=true,breaksymbol=,fontsize=\small,linenos=false]{sh}
#!/usr/bin/env sh
# Afiseaza utilizatorii care au GID-ul dat si numele grupului.
usage() {
    echo "Usage: ${0} GID" >&2
    echo "Prints users belonging to group with the specified GID." >&2
}

[ "$#" -ne 1 ] && { usage; exit 1; }

gid=$1
group_name=$(grep "^[^:]*:[^:]*:${gid}:" /etc/group | awk -F: 'NR==1{print $1}')

if [ -z "$group_name" ]; then
    echo "Nu exista grup cu GID=$gid"
    exit 1
fi

echo "Grup: $group_name (GID=$gid)"
awk -F: -v gid="$gid" '$4 == gid {print $1}' /etc/passwd
\end{minted}

\begin{center}
\includegraphics[width=.9\linewidth]{./screenshots/grep/5.png}
\end{center}
\section{Probleme de rezolvat folosind utilitarul sed}
\label{sec:org529380a}
\subsection{Se cere sa se stearga toate aparitiile cuvintelor date ca parametri din fisierul care este dat ca prim parametru.}
\label{sec:orgdf27021}
\begin{minted}[breaklines=true,breakanywhere=true,breaksymbol=,fontsize=\small,linenos=false]{sh}
#!/usr/bin/env sh
# Sterge toate aparitiile cuvintelor date ca parametri din fisierul dat.
usage() {
    echo "Usage: ${0} FILE WORD1 [WORD2 ...]" >&2
    echo "Deletes all occurrences of given words from FILE." >&2
}

[ "$#" -lt 2 ] && { usage; exit 1; }

file=$1
shift

script=''
for w in "$@"; do
    script="${script};s/${w}//g"
done

sed "$script" "$file"
\end{minted}

\begin{center}
\includegraphics[width=.9\linewidth]{./screenshots/sed/1.png}
\end{center}
\subsection{Sa se stearga toate liniile din fiserele date ca parametri care contin un anumit text indicat intr-un prim parametru.}
\label{sec:org13b2f06}
\begin{minted}[breaklines=true,breakanywhere=true,breaksymbol=,fontsize=\small,linenos=false]{sh}
#!/usr/bin/env sh
# Sterge liniile care contin textul dat din fisierele specificate.
usage() {
    echo "Usage: ${0} TEXT FILE1 [FILE2 ...]" >&2
    echo "Deletes lines containing TEXT from each file." >&2
}

[ "$#" -lt 2 ] && { usage; exit 1; }

text=$1
shift

for f in "$@"; do
    echo "===== $f ====="
    sed "/$text/d" "$f"
done
\end{minted}

\begin{center}
\includegraphics[width=.9\linewidth]{./screenshots/sed/2.png}
\end{center}
\subsection{Sa se stearga dintre primele 30 de linii din fisierele date ca parametri acelea care contin textul indicat de primul parametru.}
\label{sec:org0ce4dd9}
\begin{minted}[breaklines=true,breakanywhere=true,breaksymbol=,fontsize=\small,linenos=false]{sh}
#!/usr/bin/env sh
# Sterge dintre primele 30 linii pe cele care contin textul dat.
usage() {
    echo "Usage: ${0} TEXT FILE1 [FILE2 ...]" >&2
    echo "Deletes matching TEXT lines only within first 30 lines." >&2
}

[ "$#" -lt 2 ] && { usage; exit 1; }

text=$1
shift

for f in "$@"; do
    echo "===== $f ====="
    sed "1,30{/$text/d;}" "$f"
done
\end{minted}

\begin{center}
\includegraphics[width=.9\linewidth]{./screenshots/sed/3.png}
\end{center}
\subsection{Sa se adauge in fata fiecarei litere mici cuvintul indicat ca prim parametru. Fisierele tratate sunt ceilalti parametri din linia de comanda.}
\label{sec:org72c9457}
\begin{minted}[breaklines=true,breakanywhere=true,breaksymbol=,fontsize=\small,linenos=false]{sh}
#!/usr/bin/env sh
# Adauga in fata fiecarei litere mici cuvantul dat.
usage() {
    echo "Usage: ${0} PREFIX FILE1 [FILE2 ...]" >&2
    echo "Prefixes each lowercase letter with PREFIX in each file." >&2
}

[ "$#" -lt 2 ] && { usage; exit 1; }

prefix=$1
shift

for f in "$@"; do
    echo "===== $f ====="
    sed "s/[a-z]/${prefix}&/g" "$f"
done
\end{minted}

\begin{center}
\includegraphics[width=.9\linewidth]{./screenshots/sed/4.png}
\end{center}
\subsection{Sa se stearga toate cifrele din fisierele date ca parametri.}
\label{sec:org0f2f75b}
\begin{minted}[breaklines=true,breakanywhere=true,breaksymbol=,fontsize=\small,linenos=false]{sh}
#!/usr/bin/env sh
# Sterge toate cifrele din fisierele date.
usage() {
    echo "Usage: ${0} FILE1 [FILE2 ...]" >&2
    echo "Deletes all digits from each file." >&2
}

[ "$#" -lt 1 ] && { usage; exit 1; }

for f in "$@"; do
    echo "===== $f ====="
    sed 's/[0-9]//g' "$f"
done
\end{minted}

\begin{center}
\includegraphics[width=.9\linewidth]{./screenshots/sed/5.png}
\end{center}
\end{document}
